\chaptertitle{Crab Canon}

Tortoise: Oh, really? I guess there's nothing better for you than walking.
Achilles: And it's a perfect day for a walk. I think I'll be walking home soon.
Tortoise: That echoes my thoughts.
Achilles: So nice to run into you.
Tortoise: Why, same to you.
Achilles: Good day, Mr. T




Crab Canon                                                                        211


                 Typographical Number Theory

                 The Crab Canon and Indirect Self-Reference
THREE EXAMPLES OF indirect self-reference are found in the Crab Canon. Achilles
and the Tortoise both describe artistic creations they know-and, quite accidentally, those
creations happen to have the same structure as the Dialogue they're in. (Imagine my
surprise, when I, the author, noticed this!) Also, the Crab describes a biological structure
and that, too, has the same property. Of course, one could read the Dialogue and
understand it and somehow fail to notice that it, too, has the form of a crab canon. This
would be understanding it on one level, but not on another. To see the self-reference, one
has to look at the form, as well as the content, of the Dialogue.
        Gödel’s construction depends on describing the form, as well as the content, of
strings of the formal system we shall define in this Chapter -- Typographical Number
Theory (TNT). The unexpected twist is that, because of the subtle mapping which Gödel
discovered, the form of strings can be described in the formal system itself. Let us
acquaint ourselves with this strange system with the capacity for wrapping around.

What We Want to Be Able to Express in TNT
We'll begin by citing some typical sentences belonging to number theory; then we will
try to find a set of basic notions in terms of which all our sentences can be rephrased.
Those notions will then be given individual symbols. Incidentally, it should be stated at
the outset that the term "number theory" will refer only to properties of positive integers
and zero (and sets of such integers). These numbers are called the natural numbers.
Negative numbers play no role in this theory. Thus the word "number", when used, will
mean exclusively a natural number. And it is important -- vital-for you to keep separate in
your mind the formal system (TNT) and the rather ill-defined but comfortable old branch
of mathematics that is number theory itself; this I shall call "N".
Some typical sentences of N-number theory-are:

(1)    5 is prime.
(2)    2 is not a square.
(3)    1729 is a sum of two cubes.
(4)    No sum of two positive cubes is itself a cube.
(5)    There are infinitely many prime numbers.
(6)    6 is even.




Typographical Number Theory                                                             212

Now it may seem that we will need a symbol for each notion such as "prime” or "cube"
or "positive" -- but those notions are really not primitive. Primeness, for instance, has to
do with the factors which a number has, which in turn has to do with multiplication.
Cubeness as well is defined in terms multiplication. Let us rephrase the sentences, then,
in terms of what seem to be more elementary notions.

   (1) There do not exist numbers a and b, both greater than 1. such that 5 equals a
       times b.

   (2) There does not exist a number b, such that b times b equals 2.

   (3) There exist numbers b and c such that b times b times b, plus c times c times c,
       equals 1729.

   (4') For all numbers b and c, greater than 0, there is no number a such that a times a
        times a equals b times b times b plus c times c times c.

   (5) For each number a, there exists a number b, greater than a, with the property
       that there do not exist numbers c and d, both greater than 1, such that b equals c
       times d.

   (6') There exists a number e such that 2 times e equals 6.

This analysis has gotten us a long ways towards the basic elements of language of
number theory. It is clear that a few phrases reappear over a over:

            for all numbers b
            there exists a number b, such that
            greater than
            equals
            times
            plus
            0, 1, 2, . .

Most of these will be granted individual symbols. An exception is "greater than", which
can be further reduced. In fact, the sentence "a is greater than b" becomes

       there exists a number c, not equal to 0, such that a equals b plus c.

                                       Numerals
We will not have a distinct symbol for each natural number. Instead, we have a very
simple, uniform way of giving a compound symbol to e natural number -- very much as
we did in the pq-system. Here is notation for natural numbers:




Typographical Number Theory                                                             213

                  zero:               0
                  one:               SO
                  two:              SSO
                  three:           SSSO

                            etc.

The symbol S has an interpretation-"the successor of". Hence, the interpretation of SSO
is literally "the successor of the successor of zero". Strings of this form are called
numerals.

                                   Variables and Terms
Clearly, we need a way of referring to unspecified, or variable, numbers. For that, we will
use the letters a, b, c, d, e. But five will not be enough. We need an unlimited supply of
them, just as we had of atoms in the Propositional Calculus. We will use a similar method
for making more variables: tacking on any number of primes. (Note: Of course the
symbol "'-read "prime"-is not to be confused with prime numbers!) For instance:

                               e
                               d'
                               c"
                               b´´´...
                               a´´´´

are all variables.
         In a way it is a luxury to use the first five letters of the alphabet when we could
get away with just a and the prime. Later on, I will actually drop b, c, d, and e, which will
result in a sort of "austere" version of TNT-austere in the sense that it is a little harder to
decipher complex formulas. But for now we'll be luxurious.
         Now what about addition and multiplication? Very simple: we will use the
ordinary symbols `+' and `•'. However, we will also introduce a parenthesizing
requirement (we are now slowly slipping into the rules which define well-formed strings
of TNT). To write "b plus c" and "b times c", for instance, we use the strings

                                            (b+c)
                                            (b • c)

There is no laxness about such parentheses; to violate the convention is to produce a non-
well-formed formula. ("Formula"? I use the term instead of "string" because it is
conventional to do so. A formula is no more and no less than a string of TNT.)

Incidentally, addition and multiplication are always to be thought of as binary operations-
that is, they unite precisely two numbers, never three or more. Hence, if you wish to
translate "1 plus 2 plus 3", you have to decide which of the following two expressions
you want:


Typographical Number Theory                                                                214

                                    (SO+(SSO+SSSO))
                                    ((SO+SSO)+SSSO)

The next notion we'll symbolize is equals. That is very simple: we use ´=´.The advantage
of taking over the standard symbol used N -- nonformal number theory -- iis obvious:
easy legibility. The disadvantage is very much like the disadvantage of using the words
"point" a "line" in a formal treatment of geometry: unless one is very conscious a careful,
one may blur the distinction between the familiar meaning and strictly rule-governed
behavior of the formal symbol. In discuss geometry, I distinguished between the
everyday word and the formal to by capitalizing the formal term: thus, in elliptical
geometry, a POINT was 1 union of two ordinary points. Here, there is no such
distinction; hen mental effort is needed not to confuse a symbol with all of the association
is laden with. As I said earlier, with reference to the pq-system: the string --- is not the
number 3, but it acts isomorphically to 3, at least in the context of additions. Similar
remarks go for the string SSSO.

                        Atoms and Propositional Symbols
All the symbols of the Propositional Calculus except the letters used making atoms (P, Q,
and R) will be used in TNT, and they retain their interpretations. The role of atoms will
be played by strings which, when interpreted, are statements of equality, such as
SO=SSO or (SO • SO) Now, we have the equipment to do a fair amount of translation of
simple sentences into the notation of TNT:

            2 plus 3 equals 4:                (SSO+SSSO)=SSSSO
            2 plus 2 is not equal to 3:       ~(SSO+SSO)=SSSO
            If 1 equals 0, then 0 equals 1:   <SO=OJO=SO>

The first of these strings is an atom; the rest are compound formulas (Warning: The `and'
in the phrase "I and 1 make 2" is just another word for `plus', and must be represented by
`+' (and the requisite parentheses).)

                          Free Variables and Quantifiers
All the well-formed formulas above have the property that their interpretations are
sentences which are either true or false. There are, however, well-formed formulas which
do-not have that property, such as this one

                                       (b+SO)=SSO

Its interpretation is "b plus 1 equals 2". Since b is unspecified, there is way to assign a
truth value to the statement. It is like an out-of-context statement with a pronoun, such as
"she is clumsy". It is neither true nor false; it is waiting for you to put it into a context.
Because it is neither true nor false, such a formula is called open, and the variable b is
called a free variable.



Typographical Number Theory                                                               215

        One way of changing an open formula into a closed formula, or sentence, is by
prefixing it with a quantifier-either the phrase "there exists a number b such that , or the
phrase "for all numbers b". In the first instance, you get the sentence

       There exists a number b such that b plus 1 equals 2.

Clearly this is true. In the second instance, you get the sentence

       For all numbers b, b plus 1 equals 2.

Clearly this is false. We now introduce symbols for both of these quantifiers. These
sentences are translated into TNT-notation as follows:

                     ℑb:(b+SO)=SSO              ('ℑ' stands for `exists'.)
                     Vb:(b+SO)=SSO              ('V' stands for `all'.)

It is very important to note that these statements are no longer about unspecified
numbers; the first one is an assertion of existence, and the second one is a universal
assertion. They would mean the same thing, even if written with c instead of b:

                                     ℑc:(c+SO)=SSO `
                                     Vc:(c+SO)=SSO

A variable which is under the dominion of a quantifier is called a quantified variable. The
following two formulas illustrate the difference between free variables and quantified
variables:

                       (b.b)=SSO               (open)
                       ---ℑb:(b•b)=SSO         (closed; a sentence of TNT)

The first one expresses a property which might be possessed by some natural number. Of
course, no natural number has that property. And that is precisely what is expressed by
the second one. It is very crucial to understand this difference between a string with a free
variable, which expresses a property, and a string where the variable is quantified, which
expresses a truth or falsity. The English translation of a formula with at least one free
variable-an open formula-is called a predicate. It is a sentence without a subject (or a
sentence whose subject is an out-of-context pronoun). For instance,

                             "is a sentence without a subject"

                                  "would be an anomaly"

                      "runs backwards and forwards simultaneously"

                         "improvised a six-part fugue on demand"

are nonarithmetical predicates. They express properties which specific entities might or
might not possess. One could as well stick on a "dummy




Typographical Number Theory                                                              216

subject", such as "so-and-so". A string with free variables is like a predicate with "so-
and-so" as its subject. For instance,

                                       (SO+SO)=b

is like saying "1 plus 1 equals so-and-so". This is a predicate in the variable b. It
expresses a property which the number b might have. If one wet substitute various
numerals for b, one would get a succession of forms most of which would express
falsehoods. Here is another example of difference between open formulas and sentences:

                                  `Vb:`Vc:(b+c)=(c+b)

The above formula is a sentence representing, of course, the commutativity of addition.
On the other hand,

                                     `Vc:(b+c)=(c+b)

is an open formula, since b is free. It expresses a property which unspecified number b
might or might not have -- namely of commuting with all numbers c.

                       Translating Our Sample Sentences
This completes the vocabulary with which we will express all num theoretical statements!
It takes considerable practice to get the hang of expressing complicated statements of N
in this notation, and converse] figuring out the meaning of well-formed formulas. For this
reason return to the six sample sentences given at the beginning, and work their
translations into TNT. By the way, don't think that the translations given below are
unique-far from it. There are many -- infinitely many -- ways to express each one.
         Let us begin with the last one: "6 is even". This we rephrased in to of more
primitive notions as "There exists a number e such that 2 times e equals 6". This one is
easy
       :
                                 ℑe:(SSO. e)=SSSSSSO

Note the necessity of the quantifier; it simply would not do to write

                                   (SSO . e)=SSSSSSO

alone. This string's interpretation is of course neither true nor false; it expresses a
property which the number e might have.
        It is curious that, since we know multiplication is commutative might easily have
written

                                 ℑe:(e - SSO)=SSSSSSO

instead. Or, knowing that equality is a symmetrical relation, we might 1 chosen to write
the sides of the equation in the opposite order:


Typographical Number Theory                                                           217

                                 ℑe:SSSSSSO=(SSO • e)

Now these three translations of "6 is even" are quite different strings, and it is by no
means obvious that theoremhood of any one of them is tied to theoremhood of any of the
others. (Similarly, the fact that --p-q--- was a theorem had very little to do with the fact
that its "equivalent" string -p--q--- was a theorem. The equivalence lies in our minds,
since, as humans, we almost automatically think about interpretations, not structural
properties of formulas.)
        We can dispense with sentence 2: "2 is not a square", almost immediately:

                                     -ℑb:(b • b)=SSO

However, once again, we find an ambiguity. What if we had chosen to write it this way?

                                    Vb: -(b • b) =SSO

The first way says, "It is not the case that there exists a number b with the property that
b's square is 2", while the second way says, "For all numbers b, it is not the case that b's
square is 2." Once again, to us, they are conceptually equivalent-but to TNT, they are
distinct strings.
        Let us proceed to sentence 3: "1729 is a sum of two cubes." This one will involve
two existential quantifiers, one after the other, as follows:

              ℑb:ℑc:SSSSSS…………SSSSSO=(((b • b) • b)+((c • c) • c))

                              1729 of them

There are alternatives galore. Reverse the order of the quantifiers; switch the sides of the
equation; change the variables to d and e; reverse the addition; write the multiplications
differently; etc., etc. However, I prefer the following two translations of the sentence:

            ℑb:ℑc:(((SSSSSSSSSSO.SSSSSSSSSSO).SSSSSSSSSSO)+
     ((SSSSSSSSSO • SSSSSSSSSO) • SSSSSSSSSO))=(((b • b) • b)+((c • c) • c))

                                            and

         ℑb:ℑc:(((SSSSSSSSSSSSO.SSSSSSSSSSSSO). SSSSSSSSSSSSO)+
                     ((SO •SO) • SO))=(((b •b) •b)+((c • c) •c))

Do you see why?

                                   Tricks of the Trade

Now let us tackle the related sentence 4: "No sum of two positive cubes is itself a cube".
Suppose that we wished merely to state that 7 is not a sum of two positive cubes. The
easiest way to do this is by negating the formula


Typographical Number Theory                                                             218

which asserts that 7 is a sum of two positive cubes. This will be just like the preceding
sentence involving 1729, except that we have to add in the proviso of the cubes being
positive. We can do this with a trick: prefix variables with the symbol S, as follows:

                  ℑb:ℑc:SSSSSSSO=(((Sb • Sb) • Sb)+((Sc • Sc) -Sc))

You see, we are cubing not b and c, but their successors, which must be positive, since
the smallest value which either b or c can take on is zero. Hence the right-hand side
represents a sum of two positive cubes. In( tally, notice that the phrase "there exist
numbers b and c such that…..”) when translated, does not involve the symbol `n' which
stands for ‘and’. That symbol is used for connecting entire well-formed strings, not for
joining two quantifiers.
        Now that we have translated "7 is a sum of two positive cubes", we wish to negate
it. That simply involves prefixing the whole thing by a single (Note: you should not
negate each quantifier, even though the desired phrase runs "There do not exist numbers
b and c such that ...".) Thus we get:

                  -ℑb:ℑc:SSSSSSSO=(((Sb • Sb) • Sb)+((Sc -Sc) -Sc))

Now our original goal was to assert this property not of the number of all cubes.
Therefore, let us replace the numeral SSSSSSSO by the ((a-a)-a), which is the translation
of "a cubed":

                   ℑb:ℑc:((a •a) •a)=(((Sb •Sb) • Sb)+((Sc -Sc) -Sc))
    -
At this stage, we are in possession of an open formula, since a is still free. This formula
expresses a property which a number a might or might not have-and it is our purpose to
assert that all numbers do have that property. That is simple -- just prefix the whole thing
with a universal quantifier

                Va:-ℑb:ℑc:((a -a) • a)=(((Sb • Sb) • Sb) +((Sc -Sc) -Sc))

An equally good translation would be this:

                   --ℑa:ℑb:ℑc:((a-a) a)=(((Sb•Sb)•Sb)+((Sc•Sc)•Sc))

In austere TNT, we could use a' instead of b, and a" instead of c, and the formula would
become:

         --ℑa: ℑa': ℑa":((a • a) • a) =(((Sa' • Sa') • Sa') +((Sa" • Sa") • Sa"))

What about sentence 1: "5 is prime"? We had reworded it in this way "There do not exist
numbers a and b, both greater than 1, such equals a times b". We can slightly modify it,
as follows: "There do not exist numbers a and b such that 5 equals a plus 2, times b plus
2". This is another trick-since a and b are restricted to natural number values, this is an
adequate way to say the same thing. Now "b plus 2" could be translated into


Typographical Number Theory                                                             219

(b+SSO), but there is a shorter way to write it -- namely, SSb. Likewise, "c plus 2" can
be written SSc. Now, our translation is extremely concise:

                               ℑb: ℑc:SSSSSO=(SSb • SSc)

Without the initial tilde, it would be an assertion that two natural numbers do exist,
which, when augmented by 2, have a product equal to 5. With the tilde in front, that
whole statement is denied, resulting in an assertion that 5 is prime.
       If we wanted to assert that d plus e plus 1, rather than 5, is prime, the most
economical way would be to replace the numeral for 5 by the string (d+Se):

                                ℑb: ℑc:(d+Se)=(SSb SSc)

Once again, an open formula, one whose interpretation is neither a true nor a false
sentence, but just an assertion about two unspecified numbers, d and e. Notice that the
number represented by the string (d+Se) is necessarily greater than d, since one has
added to d an unspecified but definitely positive amount. Therefore, if we existentially
quantify over the variable e, we will have a formula which asserts that:

There exists a number which is greater than d and which is prime.

                             ℑe:- ℑb:3c:(d+Se)=(SSb • SSc)

Well, all we have left to do now is to assert that this property actually obtains, no matter
what d is. The way to do that is to universally quantify over the variable d:

                            Vd:3e:-3b:3c:(d+Se)=(SSb •SSc)

That's the translation of sentence 5!


                            Translation Puzzles for You
This completes the exercise of translating all six typical number-theoretical sentences.
However, it does not necessarily make you an expert in the notation of TNT. There are
still some tricky issues to be mastered. The following six well-formed formulas will test
your understanding of TNT notation. What do they mean? Which ones are true (under
interpretation, of course), and which ones are false? (Hint: the way to tackle this exercise
is to move leftwards. First, translate the atom; next, figure out what adding a single
quantifier or a tilde does; then move leftwards, adding another quantifier or tilde; then
move leftwards again, and do the same.)

                                   -Vc: ℑb:(SSO • b)=c

                                   Vc:- ℑb:(SSO • b)=c



Typographical Number Theory                                                             220

                                   Vc: ℑb:---(SSO • b)=c
                                   ~ℑb:Vc:(SSO • b)=c
                                   ℑb:- Vc:(SSO • b)=c
                                   ℑb: Vc:-(SSO • b)=c

(Second hint: Either four of them are true and two false, or four false and two true.)

                       How to Distinguish True from False?
At this juncture, it is worthwhile pausing for breath and contempt what it would mean to
have a formal system that could sift out the true from the false ones. This system would
treat all these strings-which look like statements-as designs having form, but no content.
An( system would be like a sieve through which could pass only designs v special style-
the "style of truth". If you yourself have gone through ti formulas above, and have
separated the true from the false by this about meaning, you will appreciate the subtlety
that any system would to have, that could do the same thing-but typographically! The
bout separating the set of true statements from the set of false statements written in the
TNT-notation) is anything but straight; it is a boundary with many treacherous curves
(recall Fig. 18), a boundary of which mathematicians have delineated stretches, here and
there, working over hundreds years. Just think what a coup it would be to have a
typographical m( which was guaranteed to place any formula on the proper side o border!

                           The Rules of Well-Formedness
It is useful to have a table of Rules of Formation for well-formed formulas This is
provided below. There are some preliminary stages, defining numerals, variables, and
terms. Those three classes of strings are ingredients of well-formed formulas, but are not
in themselves well-formed. The smallest well-formed formulas are the atoms; then there
are ways of compounding atoms. Many of these rules are recursive lengthening rules, in
that they take as input an item of a given class and produce a longer item of the class. In
this table, I use `x' and 'y' to stand for well-formed formulas, and `s', `t', and `u' to stand
for other kinds of TNT-strings. Needless to say, none of these five symbols is itself a
symbol of TNT.

NUMERALS.
    0 is a numeral.
    A numeral preceded by S is also a numeral.
    Examples: 0 SO S50 SSSO SSSSO SSSSSO




Typographical Number Theory                                                                221

VARIABLES.
         a is a variable. If we're not being austere, so are b, c, d and e. A variable followed
         by a prime is also a variable.
         Examples: a b' c" d"' a""
TERMS.
         All numerals and variables are terms.
         A term preceded by S is also a term.
         If s and t are terms, then so are (s+ t) and (s • t).
         Examples: 0 b SSa' (SO • (SSO+c)) S(Sa • (Sb • Sc))

         TERMS may be divided into two categories:

                 (1) DEFINITE terms. These contain no variables.
                     Examples: 0 (SO+SO) SS((SSO.SSO)+(SO.SO))
                 (2) INDEFINITE terms. These contain variables.
                    Examples: b Sa (b+SO) (((SO+SO)+SO)+e)
        The above rules tell how to make parts of well-formed formulas; the remaining
rules tell how to make complete well-formed formulas.
ATOMS.
         If s and t are terms, then s = t is an atom.
         Examples: SO=0 (SS0+SS0)=5SSS0 5(b+c)=((c•d).e)
         If an atom contains a variable u, then u is free in it. Thus there are
         four free variables in the last example.
NEGATIONS.
         A well-formed formula preceded by a tilde is well-formed.
         Examples: ~S0=0 ~ℑb:(b+b)=SO -<O=0⊃S0=O> ~b=SO
         The quantification status of a variable (which says whether the variable is
                free or quantified) does not change under negation.
COMPOUNDS.
         If x and y are well-formed formulas, and provided that no variable which is free in
         one is quantified in the other, then the following are all well-formed formulas:
         < x∧ y>, < x∨ y>, < x⊃ y>.
         Examples: <O=O∧~-0=0> <b=b∨~ℑc:c=b>
         <SO=O⊃Vc:~ℑb:(b+b)=c>
         The quantification status of a variable doesn't change here.
QUANTI FI CATIONS.
       If u is a variable, and x is a well-formed formula in which u is free then the
following strings are well-formed formulas:
         ℑu: x and Vu: x.
         Examples: Vb:<b=b∨~ℑc:c=b>           ∨c:~ℑb:(b+b)=c ~ℑc:Sc=d
OPEN FORMULAS contain at least one free variable.
         Examples: --c=c b=b <Vb:b=bn---c=c>
CLOSED FORMULAS (SENTENCES) contain no free variables.
         Examples: 50=0     ~Vd:d=0      ℑc:<Vb:b=b∧~c=c>




Typographical Number Theory                                                                222

This completes the table of Rules of Formation for the well-formed formulas of TNT.

                        A Few More Translation Exercises
And now, a few practice exercises for you, to test your understanding of the notation of
TNT. Try to translate the first four of the following N-sentences into TNT-sentences, and
the last one into an open formed formula.

                                All natural numbers are equal to 4.

                     There is no natural number which equals its own square.

                       Different natural numbers have different successors.

                              If 1 equals 0, then every number is odd.

                                         b is a power of 2.

The last one you may find a little tricky. But it is nothing, compared to this one:

       b is a power of 10.

Strangely, this one takes great cleverness to render in our notation. I would caution you to
try it only if you are willing to spend hours and hours on it -- and if you know quite a bit
of number theory!

                             A Non typographical System
This concludes the exposition of the notation of TNT; however, we still left with the
problem of making TNT into the ambitious system which we have described. Success
would justify the interpretations which we given to the various symbols. Until we have
done that, however, particular interpretations are no more justified than the "horse-apple
happy" interpretations were for the pq-system's symbols.
        Someone might suggest the following way of constructing TNT: (1|) Do not have
any rules of inference; they are unnecessary, because (2) We take as axioms all true
statements of number theory (as written in TNT-notation). What a simple prescription!
Unfortunately it is as empty as instantaneous reaction says it is. Part (2) is, of course, not
a typographical description of strings. The whole purpose of TNT is to figure out if and
how it is possible to- characterize the true strings typographically.

                    The Five Axioms and First Rules of TNT
Thus we will follow a more difficult route than the suggestion above; we will have
axioms and rules of inference. Firstly, as was promised, all of the rules of the
Propositional Calculus are taken over into TNT. Therefore one theorem of TNT will be
this one:


                                      <S0=0∨~S0=0>


Typographical Number Theory                                                               223

which can be derived in the same way as <P∨-P> was derived.
       Before we give more rules, let us give the five axioms of TNT:

Axiom 1: Va:~Sa=O

Axiom 2: Va:(a+O)=a

Axiom 3: Va:Vb:(a+Sb)=S(a+b)

Axiom 4: Va:(a-O)=O

Axiom 5: Va:Vb:(a-Sb)=((a-b)+a)

(In the austere versions, use a' instead of b.) All of them are very simple to understand.
Axiom 1 states a special fact about the number 0; Axioms 2 and 3 are concerned with the
nature of addition; Axioms 4 and 5 are concerned with the nature of multiplication, and in
particular with its relation to addition.

                             The Five Peano Postulates
By the way, the interpretation of Axiom 1-"Zero is not the successor of any natural
number"-is one of five famous properties of natural numbers first explicitly recognized
by the mathematician and logician Giuseppe Peano, in 1889. In setting out his postulates,
Peano was following the path of Euclid in this way: he made no attempt to formalize the
principles of reasoning, but tried to give a small set of properties of natural numbers from
which everything else could be derived by reasoning. Peano's attempt might thus be
considered "semiformal". Peano's work had a significant influence, and thus it would be
good to show Peano's five postulates. Since the notion of "natural number" is the one
which Peano was attempting to define, we will not use the familiar term "natural
number", which is laden with connotation. We will replace it with the undefined term
djinn, a word which comes fresh and free of connotations to our mind. Then Peano's five
postulates place five restrictions on djinns. There are two other undefined terms: Genie,
and meta. I will let you figure out for yourself what usual concept each of them is
supposed to represent. The five Peano postulates:

   (1) Genie is a djinn.
   (2) Every djinn has a mesa (which is also a djinn).
   (3) Genie is not the mesa of any djinn. (4) Different djinns have different metas.
   (5) If Genie has X, and each djinn relays X to its mesa, then all djinns get X.

In light of the lamps of the Little Harmonic Labyrinth, we should name the set of all
djinns "GOD". This harks back to a celebrated statement by the German mathematician
and logician Leopold Kronecker, archenemy of Georg Cantor: "God made the natural
numbers; all the rest is the work of man."




Typographical Number Theory                                                             224

         You may recognize Peano's fifth postulate as the principle of mathematical
induction-another term for a hereditary argument. Peano he that his five restrictions on
the concepts "Genie", "djinn", and "mesa" so strong that if two different people formed
images in their minds o concepts, the two images would have completely isomorphic
structures. example, everybody's image would include an infinite number of distinct
djinns. And presumably everybody would agree that no djinn coins with its own meta, or
its meta's meta, etc.
         Peano hoped to have pinned down the essence of natural numbers in his five
postulates. Mathematicians generally grant that he succeeded that does not lessen the
importance of the question, "How is a true statement about natural numbers to be
distinguished from a false one?" At answer this question, mathematicians turned to totally
formal systems, as TNT. However, you will see the influence of Peano in TNT, because
all of his postulates are incorporated in TNT in one way or another.

            New Rules of TNT: Specification and Generalization
Now we come to the new rules of TNT. Many of these rules will allow reach in and
change the internal structure of the atoms of TNT. In sense they deal with more
"microscopic" properties of strings than the of the Propositional Calculus, which treat
atoms as indivisible units. example, it would be nice if we could extract the string -SO=O
from the first axiom. To do this we would need a rule which permits us to di universal
quantifier, and at the same time to change the internal strut of the string which remains, if
we wish. Here is such a rule:

RULE OF SPECIFICATION: Suppose u is a variable which occurs inside string x. If the
string Vu:x is a theorem, then so is x, and so an strings made from x by replacing u,
wherever it occurs, by one the same term.

       (Restriction: The term which replaces u must not contain any vat that is quantified
       in x.)

The rule of specification allows the desired string to be extracted Axiom 1. It is a one-
step derivation:

                          Va -Sa=0                       axiom 1
                          ~S0=0                          specification

Notice that the rule of specification will allow some formulas which co: free variables
(i.e., open formulas) to become theorems. For example following strings could also be
derived from Axiom 1, by specification:

                                           Sa=0
                                       ~S(c+SSO)=0

There is another rule, the rule of generalization, which allows us to put


Typographical Number Theory                                                              225

back the universal quantifier on theorems which contain variables that became free as a
result of usage of specification. Acting on the lower string, for example, generalization
would give:

                                   Vc:~S(c+SSO)=O

Generalization undoes the action of specification, and vice versa. Usually, generalization
is applied after several intermediate steps have transformed the open formula in various
ways. Here is the exact statement of the rule:

RULE OF GENERALIZATION: Suppose x is a theorem in which u, a variable, occurs
free. Then Vu:x is a theorem.

       ( Restriction: No generalization is allowed in a fantasy on any variable which
       appeared free in the fantasy's premise.)

The need for restrictions on these two rules will shortly be demonstrated explicitly.
Incidentally, this generalization is the same generalization as was mentioned in Chapter
II, in Euclid's proof about the infinitude of primes. Already we can see how the symbol-
manipulating rules are starting to approximate the kind of reasoning which a
mathematician uses.

                            The Existential Quantifier
These past two rules told how to take off universal quantifiers and put them back on; the
next two rules tell how to handle existential quantifiers.

RULE OF INTERCHANGE: Suppose u is a variable. Then the strings Vu:- and -3u: are
interchangeable anywhere inside any theorem.

For example, let us apply this rule to Axiom 1:

                                Va:-Sa=O   axiom 1
                               ~ℑa:Sa=O  interchange

By the way, you might notice that both these strings are perfectly natural renditions, in
TNT, of the sentence "Zero is not the successor of any natural number". Therefore it is
good that they can be turned into each other with ease.
       The next rule is, if anything, even more intuitive. It corresponds to the very
simple kind of inference we make when we go from "2 is prime" to "There exists a
prime". The name of this rule is self-explanatory:

RULE OF EXISTENCE: Suppose a term (which may contain variables as long as they
are free) appears once, or multiply, in a theorem. Then any (or several, or all) of the
appearances of the term may be replaced by a variable which otherwise does not occur in
the theorem, and the corresponding existential quantifier must be placed in front.

Let us apply the rule to --as usual--Axiom 1:



Typographical Number Theory                                                           226

                                   Va:-Sa=O     axiom 1
                                   ℑb:Va:-Sa=b existence

You might now try to shunt symbols, according to rules so far giver produce the theorem
~Vb: ℑa:Sa=b.

                      Rules of Equality and Successorship
We have given rules for manipulating quantifiers, but so far none for symbols `=' and 'S'.
We rectify that situation now. In what follows, r, s, t all stand for arbitrary terms.
RULES OF EQUALITY:

       SYMMETRY: If r = s is a theorem, then so is s = r.
       TRANSITIVITY: If r = s and s = t are theorems, then so is r = t.

RULFS OF SUCCESSORSHIP:

       ADD S: If r = t is a theorem, then Sr = St is a theorem.

       DROP S: If Sr = St is a theorem, then r = t is a theorem.

Now we are equipped with rules that can give us a fantastic variet theorems. For
example, the following derivations yield theorems which pretty fundamental:

(1)    Va: Vb:(a+Sb)=S(a+b)                         axiom 3
(2)    Vb:(SO+Sb)=S(SO+b)                           specification (SO for a)
(3)    (SO+SO)=S(SO+O)                              specification (0 for b)
(4)    Va:(a+O)=a                                   axiom 2
(5)    (SO+O)=SO                                    specification (SO for a)
(6)    S(SO+O)=SSO                                  add S
(7)    (SO+SO)=SSO                                  transitivity (lines 3,6)

                               *      *    *    *      *

(1)    Va: Vb:(a-Sb)=((a-b)+a)                      axiom 5
(2)    Vb:(SO•Sb)=((SO•b)+SO)                       specification (SO for a)
(3)    (SO.SO)=((SO.O)+SO)                          specification (0 for b)
(4)    Va: Vb:(a+Sb)=S(a+b)                                 axiom 3
(5)    Vb:((SO.O)+Sb)=S((5O O)+b)                   specification ((SO-0) for a)
(6)    ((SO .0)+SO)=S((SO.0)+0)                     specification (0 for b)
(7)    Va:(a+O)=a                                   axiom 2
(8)    ((SO.O)+0)=(SO.O)                            specification ((S0.0) for a)
(9)    Va:(a.0)=0                                   axiom 4
(10)   (S0-0)=0                                     specification (SO for a)
(11)   ((SO.O)+O)=O                                 transitivity (lines 8,10)
(12)   S((SO.0)+0)=SO                               add S
(13)   ((SO -0)+SO)=SO                              transitivity (lines 6,12)
(14)   (SO.SO)=SO                                   transitivity (lines 3,13)


Typographical Number Theory                                                           227

                                     Illegal Shortcuts
Now here is an interesting question: "How can we make a derivation for the string 0=0?"
It seems that the obvious route to go would be first to derive the string Va:a=a, and then
to use specification. So, what about the following "derivation" of Va:a=a ... What is
wrong with it? Can you fix it up?

       (1)     Va:(a+0)=a              axiom 2
       (2)     Va:a=(a+0)              symmetry
       (3)     Va:a=a                  transitivity (lines 2,1)

I gave this mini-exercise to point out one simple fact: that one should not jump too fast in
manipulating symbols (such as `=') which are familiar. One must follow the rules, and not
one's knowledge of the passive meanings of the symbols. Of course, this latter type of
knowledge is invaluable in guiding the route of a derivation.

             Why Specification and Generalization Are Restricted
Now let us see why there are restrictions necessary on both specification and
generalization. Here are two derivations. In each of them, one of the restrictions is
violated. Look at the disastrous results they produce:

(1)    [                               push
(2)        a=0                         premise
(3)        Va:a=0                      generalization (Wrong!)
(4)        Sa=O                        specification
(5)    ]                               pop
(6)    <a=O⊃Sa=O>                              fantasy rule
(7)    Va:<a=O⊃Sa=O>                   generalization
(8)    <O=O⊃SO=0>                      specification
(9)    0=0                             previous theorem
(10)   S0=0                            detachment (lines 9,8)

This is the first disaster. The other one is via faulty specification.

(1)    Va:a=a                          previous theorem
(2)    Sa=Sa                           specification
(3)    ℑb:b=Sa                         existence
(4)    Va: ℑb:b=Sa                     generalization
(5)    ℑb:b=Sb                         specification (Wrong!)

So now you can see why those restrictions are needed.
       Here is a simple puzzle: translate (if you have not already done so) Peano's fourth
postulate into TNT-notation, and then derive that string as a theorem.




Typographical Number Theory                                                             228

                                 Something Is Missing
Now if you experiment around for a while with the rules and axioms of TNT so far
presented, you will find that you can produce the following pyramidal family of theorems
(a set of strings all cast from an identical mold, differing from one another only in that the
numerals 0, SO, SSO, and s have been stuffed in):

                                        (0+0)=0
                                      (O+SO)=S0
                                     (O+SSO)=SSO
                                    (O+SSSO)=SSSO
                                   (O+SSSSO)=SSSSO

                                             etc.

As a matter of fact, each of the theorems in this family can be derived the one directly
above it, in only a couple of lines. Thus it is a so "cascade" of theorems, each one
triggering the next. (These theorem very reminiscent of the pq-theorems, where the
middle and right-] groups of hyphens grew simultaneously.)
         Now there is one string which we can easily write down, and v summarizes the
passive meaning of them all, taken together. That un sally quantified summarizing string
is this:

                                        Va:(O+a)=a

Yet with the rules so far given, this string eludes production. Ti produce it yourself if you
don't believe me.
        You may think that we should immediately remedy the situation the following

(PROPOSED) RULE OF ALL: If all the strings in a pyramidal family are theorems, then
so is the universally quantified string which summarizes them.

The problem with this rule is that it cannot be used in the M-mode. people who are
thinking about the system can ever know that an infinite set of strings are all theorems.
Thus this is not a rule that can be stuck i any formal system.

                ω-Incomplete Systems and Undecidable Strings

So we find ourselves in a strange situation, in which we can typographically produce
theorems about the addition of any specific numbers, but even a simple string as the one
above, which expresses a property of addition in general, is not a theorem. You might
think that is not all that strange, we were in precisely that situation with the pq-system.
However, the pq-system had no pretensions about what it ought to be able to do; and ii
fact




Typographical Number Theory                                                               229

there was no way to express general statements about addition in its symbolism, let alone
prove them. The equipment simply was not there, and it did not even occur to us to think
that the system was defective. Here, however, the expressive capability is far stronger,
and we have correspondingly higher expectations of TNT than of the pq-system. If the
string above is not a theorem, then we will have good reason to consider TNT to be
defective. As a matter of fact, there is a name for systems with this kind of defect-they
are called ω-incomplete. (The prefix 'ω'-'omega'- comes from the fact that the totality of
natural numbers is sometimes denoted by `ω'.) Here is the exact definition:

   A system is ω-incomplete if all the strings in a pyramidal family are theorems, but
   the universally quantified summarizing string is not a theorem.

   Incidentally, the negation of the above summarizing string

                                     ~Va:(O+a)=a

-is also a nontheorem of TNT. This means that the original string is undecidable within
the system. If one or the other were a theorem, then we would say that it was decidable.
Although it may sound like a mystical term, there is nothing mystical about
undecidability within a given system. It is only a sign that the system could be extended.
For example, within absolute geometry, Euclid's fifth postulate is undecidable. It has to
be added as an extra postulate of geometry, to yield Euclidean geometry; or conversely,
its negation can be added, to yield non-Euclidean geometry. If you think back to
geometry, you will remember why this curious thing happens. It is because the four
postulates of absolute geometry simply do not pin down the meanings of the terms
"point" and "line", and there is room for different extensions of the notions. The points
and lines of Euclidean geometry provide one kind of extension of the notions of "point"
and "line"; the POINTS and LINES of non-Euclidean geometry, another. However, using
the pre-flavored words "point" and "line" tended, for two millennia, to make people
believe that those words were necessarily univalent, capable of only one meaning.

                                Non-Euclidean TNT
We are now faced with a similar situation, involving TNT. We have adopted a notation
which prejudices us in certain ways. For instance, usage of the symbol `+'tends to make
us think that every theorem with a plus sign in it ought to say something known and
familiar and "sensible" about the known and familiar operation we call "addition".
Therefore it would run against the grain to propose adding the following "sixth axiom":

                                      ~Va:(0+a)=a




Typographical Number Theory                                                           230

It doesn't jibe with what we believe about addition. But it is one possible extension of
TNT, as we have so far formulated TNT. The system which uses this as its sixth axiom is
a consistent system, in the sense of not has, two theorems of the form x and - x. However,
when you juxtapose this "sixth axiom" with the pyramidal family of theorems shown
above, you will probably be bothered by a seeming inconsistency between the family and
the new axiom. But this kind of inconsistency is riot so damaging as the other kind
(where x and x are both theorems). In fact, it is not a true inconsistency, because there is
a way of interpreting the symbols so that everything comes out all right.

                   ω-Inconsistency Is Not the Same as Inconsistency

This kind of inconsistency, created by the opposition of (1) a pyramidal family of
theorems which collectively assert that all natural numbers have some property, and (2) a
single theorem which seems to assert that not all numbers have it, is given the name of w-
inconsistency. An w-inconsistent system is more like the at-the-outset-distasteful-but-in-
the-end-accept non-Euclidean geometry. In order to form a mental model of what is
going on, you have to imagine that there are some "extra", unsuspected numbers--let us
not call them "natural", but supernatural numbers-which have no numerals. Therefore,
facts about them cannot be represented in the pyramidal family. (This is a little bit like
Achilles' conception GOD-as a sort of "superdjinn", a being greater than any of the djinn
This was scoffed at by the Genie, but it is a reasonable image, and may I you to imagine
supernatural numbers.)
         What this tells us is that the axioms and rules of TNT, as so presented, do not
fully pin down the interpretations for the symbol TNT. There is still room for variation in
one's mental model of the notions they stand for. Each of the various possible extensions
would pin d, some of the notions further; but in different ways. Which symbols we begin
to take on "distasteful" passive meanings, if we added the "s axiom" given above? Would
all of the symbols become tainted, or we some of them still mean what we want them to
mean? I will let you tt about that. We will encounter a similar question in Chapter XIV,
discuss the matter then. In any case, we will not follow this extension r but instead go on
to try to repair the w-incompleteness of TNT.

                                     The Last Rule

The problem with the "Rule of All" was that it required knowing that all lines of an
infinite pyramidal family are theorems -- too much for a finite being. But suppose that
each line of the pyramid can be derived from its predecessor in a patterned way. Then
there would be a finite reason accounting for the fact that all the strings in the pyramid
are theorems. The trick then, is to find the pattern that causes the cascade, and show that




Typographical Number Theory                                                             231

pattern is a theorem in itself. That is like proving that each djinn passes a message to its
meta, as in the children's game of "Telephone". The other thing left to show is that Genie
starts the cascading message-that is, to establish that the first line of the pyramid is a
theorem. Then you know that GOD will get the message!
        In the particular pyramid we were looking at, there is a pattern, captured by lines
4-9 of the derivation below.

   (1) Va:Vb:(a+Sb)=S(a+b)                   axiom 3
   (2) Vb:(O+Sb)=S(O+b)                      specification
   (3) (O+Sb)=S(O+b)                         specification
   (4) [                                     push
   (5) (O+b)=b                               premise
   (6) S(O+b)=Sb                             add S
   (7) (O+Sb)=S(O+b)                         carry over line 3
   (8) (O+Sb)=Sb                             transitivity
   (9) ]                                     pop

The premise is (O+b)=b; the outcome is (O+Sb)=Sb.
        The first line of the pyramid is also a theorem; it follows directly from Axiom 2.
All we need now is a rule which lets us deduce that the string which summarizes the
entire pyramid is itself a theorem. Such a rule will he a formalized statement of the fifth
Peano postulate.
        To express that rule, we need a little notation. Let us abbreviate a well-formed
formula in which the variable a is free by the following notation:

                                           X{a}

(There may be other free variables, too, but that is irrelevant.) Then the notation X{Sa/a}
will stand for that string but with every occurrence of a replaced by Sa. Likewise, X{0/a}
would stand for the same string, with each appearance of a replaced by 0.
        A specific example would be to let X{a} stand for the string in question: (O+a)=a.
Then X{Sa/a} would represent the string (O+Sa)=Sa, and X{0/a} would represent
(0+0)=0. (Warning: This notation is not part of TNT; it is for our convenience in talking
about TNT.)
        With this new notation, we can state the last rule of TNT quite precisely:

RULE OF INDUCTION: Suppose u is a variable, and X{u} is a well-formed formula in
     which u occurs free. If both Vu:< X{u}⊃ X{Su/u}> and X{0/u} are theorems,
     then Vu: X{u} is also a theorem.

This is about as close as we can come to putting Peano's fifth postulate into TNT. Now
let us use it to show that Va:(O+a)=a is indeed a theorem in TNT. Emerging from the
fantasy in our derivation above, we can apply the fantasy rule, to give us

(10)   <(O+b)=b⊃(O+Sb)=Sb>                   fantasy rule
(11)   Vb:<(O+b)=b⊃(O+Sb)=Sb>                generalization


Typographical Number Theory                                                             232

This is the first of the two input theorems required by the induction The other
requirement is the first line of the pyramid, which we have. Therefore, we can apply the
rule of induction, to deduce what we wanted.

                                     `Vb:(O+b)=b

Specification and generalization will allow us to change the variable from b to a; thus
Va:(O+a)=a is no longer an undecidable string of TNT..

                                  A Long Derivation

Now I wish to present one longer derivation in TNT, so that you ca what one is like, and
also because it proves a significant, if simple, fact of number theory.

(1) Va: Vb:(a+Sb)=S(a+b)                                  axiom 3
(2) Vb:(d+Sb)=S(d+b)                                      specification
(3) (d+SSc)=S(d+Sc)                                       specificatic
(4) b:(Sd+Sb)=S(Sd+b)                                     specification (line 1)
(5) (Sd+Sc)-S(Sd+c)                                       specification
 6) S(Sd+c)=(Sd+Sc)                                       symmetry
(7) [                                                     push
(8)    Vd:(d+Sc)=(Sd+c)                                   premise
(9)    (d+Sc)=(Sd+c)                                      specification
(10) S(d+Sc)=S(Sd+c)                                      add S
(11) (d+SSc)=S(d+Sc)                                      carry over 3
(12) (d+SSc)=S(Sd+c)                                      transitivity
(13) S(Sd+c)=(Sd+Sc)                                      carry over 6
(14) (d+SSc)=(Sd+Sc)                                      transitivity
(15) Vd:(d+SSc)=(Sd+Sc)                                   generalization
(16) ]                                                    pop
(17) <Vd:(d+5c)=(Sd+c)⊃Vd:(d+SSc)=(Sd+Sc)>                fantasy rule
(18) Vc:< Vd:(d+Sc)=(Sd+c) ⊃Vd:(d+SSc)=(Sd+Sc)>           generalization

                                    *   *   *   *   *
(19) (d+S0)=5(d+0)                                        specification (line 2)
(20) Va:(a+0)=a                                           axiom 1
(21) (d+0)=d                                              specification
(22) S(d+0)=Sd                                            add S
(23) (d+SO)=Sd                                            transitivity (lines 19,2)
(24) (Sd+O)=Sd                                            specification (line 20)
(25) Sd=(Sd+O)                                            symmetry




Typographical Number Theory                                                           233

(26) (d+SO)=(Sd+o)                                       transitivity (lines 23,25)
(27) Vd:(d+5O)=(Sd+O)                                    generalization

                                 *   *   *   *   *

(28) Vc: Vd:(d+Sc)=(Sd+c)                                induction (lines 18,27)

                  [S can be slipped back and forth in an addition]
                                  * * * * *

(29) Vb:(c+Sb)=S(c+b)                                    specification (line 1)
(30) (c+Sd)=S(c+d)                                       specification
(31) Vb:(d+Sb)=S(d+b)                                    specification (line 1)
(32) (d+Sc)=S(d+c)                                       specification
(33) S(d+c)=(d+Sc)                                       symmetry
(34) bed:(d+Sc)=(Sd+c)                                   specification (line 28)
(35) (d+Sc)=(Sd+c)                                       specification
(36) [                                                   push
(37) Vc:(c+d)=(d+c)                                      premise
(38) (c+d)=(d+c)                                         specification
(39) S(c+d)=S(d+c)                                       add S
(40) (c+Sd)=S(c+d)                                       carry over 30
(41) (c+Sd)=S(d+c)                                       transitivity
(42) S(d+c)=(d+Sc)                                       carry over 33
(43) (c+Sd)=(d+Sc)                                       transitivity
(44) (d+Sc)=(Sd+c)                                       carry over 35
(45) (c+Sd)=(Sd+c)                                       transitivity
(46) Vc:(c+Sd)=(Sd+c)                                    generalization
(47) ]                                                   pop
(48) <Ve:(c+d)=(d+c) ⊃Vc:(c+Sd)=(Sd+c)>                  fantasy rule
(49) Vd:< Vc:(c+d)=(d+c) ⊃Vc:(c+Sd)=(Sd+c)>              generalization

                  [If d commutes with every c, then Sd does too.
                                * * * * *
(50) (c+O)=c                                           specification (line 20)
(51) Va:(O+a)=a                                        previous theorem
(52) (O+c)=c                                           specification
(53) c=(O+c)                                           symmetry




Typographical Number Theory                                                           234

(54) (c+0)=(0+c)                                               transitivity (lines 50,53)
(55) Vc:(c+0)=(O+c)                                            generalization

                                 [0 commutes with every c.]

(56) Vd: Vc:(c+d)=(d+c)                                        induction (lines 49,55)

                       [Therefore, every d commutes with every c.]

                          Tension and Resolution in TNT
TNT has proven the commutativity of addition. Even if you do not follow this derivation
in detail, it is important to realize that, like a piece of music, it has its own natural
"rhythm". It is not just a random walk that happens to have landed on the desired last
line. I have inserted "breathing marks” to show some of the "phrasing" of this derivation.
Line 28 in particular turning point in the derivation, something like the halfway point it
AABB type of piece, where you resolve momentarily, even if not in the t key. Such
important intermediate stages are often called "lemmas".
         It is easy to imagine a reader starting at line 1 of this derivation ignorant of where
it is to end up, and getting a sense of where it is going as he sees each new line. This
would set up an inner tension, very much the tension in a piece of music caused by chord
progressions that let know what the tonality is, without resolving. Arrival at line 28 w,
confirm the reader's intuition and give him a momentary feeling of satisfaction while at
the same time strengthening his drive to progress tow what he presumes is the true goal.
         Now line 49 is a critically important tension-increaser, because of "almost-there"
feeling which it induces. It would be extremely unsatisfactory to leave off there! From
there on, it is almost predictable how things must go. But you wouldn't want a piece of
music to quit on you just when had made the mode of resolution apparent. You don't
want to imagine ending-you want to hear the ending. Likewise here, we have to c things
through. Line 55 is inevitable, and sets up all the final tension which are resolved by Line
56.
         This is typical of the structure not only of formal derivations, but of informal
proofs. The mathematician's sense of tension is intimately related to his sense of beauty,
and is what makes mathematics worthy doing. Notice, however, that in TNT itself, there
seems to be no reflection of these tensions. In other words, TNT doesn't formalize the
notions of tension and resolution, goal and subgoal, "naturalness" and "inevitable any
more than a piece of music is a book about harmony and rhythm. Could one devise a
much fancier typographical system which is aware of the tensions and goals inside
derivations?




Typographical Number Theory                                                                 235

                   Formal Reasoning vs. Informal Reasoning
         I would have preferred to show how to derive Euclid's Theorem (the infinitude of
primes) in TNT, but it would probably have doubled the length of the book. Now after
this theorem, the natural direction to go would be to prove the associativity of addition,
the commutativity and associativity of multiplication and the distributivity of
multiplication over addition. These would give a powerful base to work from.
         As it is now formulated, TNT has reached "critical mass" (perhaps a strange
metaphor to apply to something called "TNT"). It is of the same strength as the system of
Principia Mathematica; in TNT one can now prove every theorem which you would find
in a standard treatise on number theory. Of course, no one would claim that deriving
theorems in TNT is the best way to do number theory. Anybody who felt that way would
fall in the same class of people as those who think that the best way to know what 1000 x
1000 is, is to draw a 1000 by 1000 grid, and count all the squares in it ... No; after total
formalization, the only way to go is towards relaxation of the formal system. Otherwise,
it is so enormously unwieldy as to be, for all practical purposes, useless. Thus, it is
important to embed TNT within a wider context, a context which enables new rules of
inference to be derived, so that derivations can be speeded up. This would require
formalization of the language in which rules of inference are expressed-that is, the
metalanguage. And one could go considerably further. However, none of these speeding-
up tricks would make TNT any more powerful; they would simply make it more usable.
The simple fact is that we have put into TNT every mode of thought that number
theorists rely on. Embedding it in ever larger contexts will not enlarge the space of
theorems; it will just make working in TNT-or in each "new, improved version"-look
more like doing conventional number theory.

                          Number Theorists Go out of Business
        Suppose that you didn't have advance knowledge that TNT will turn out to be
incomplete, but rather, expected that it is complete-that is, that every true statement
expressible in the TNT-notation is a theorem. In that case, you could make a decision
procedure for all of number theory. The method would be easy: if you want to know if N-
statement X is true or false, code it into TNT-sentence x. Now if X is true, completeness
says that x is a theorem; and conversely, if not-X is true, then completeness says that ~x
is a theorem. So either x or ~x must be a theorem, since either X or not-X is true. Now
begin systematically enumerating all the theorems of TNT, in the way we did for the
MIU-system and pq-system. You must come to x or ~x after a while; and whichever one
you hit tells you which of X and not-X is true. (Did you follow this argument? It crucially
depends on your being able to hold separate in your mind the formal system TNT and its
informal counterpart N. Make sure you understand it.) Thus, in prince-




Typographical Number Theory                                                             236

ple, if TNT were complete, number theorists would be put out of business any question
in their field could be resolved, with sufficient time, in a purely mechanical way. As it
turns out, this is impossible, which, depending on your point of view, is a cause either for
rejoicing, or for mourning.


                                      Hilbert's Program
        The final question which we will take up in this Chapter is whether should have
as much faith in the consistency of TNT as we did consistency of the Propositional
Calculus; and, if we don't, whether possible to increase our faith in TNT, by proving it to
be consistent could make the same opening statement on the "obviousness" of TNT’s
consistency as Imprudence did in regard to the Propositional Calculus namely, that each
rule embodies a reasoning principle which we believe in, and therefore to question the
consistency of TNT is to question our own sanity. To some extent, this argument still
carries weight-but not quite so much weight as before. There are just too many rules of
inference and some of them just might be slightly "off ". Furthermore, how do we know
that this mental model we have of some abstract entities called "natural numbers" is
actually a coherent construct? Perhaps our own thought processes, those informal
processes which we have tried to capture in the formal rules of the system, are themselves
inconsistent! It is of course not the kind of thing we expect, but it gets more and more
conceivable that our thoughts might lead us astray, the more complex the subject matter
gets-and natural numbers are by no means a trivial subject matter. Prudence's cry for a
proof of consistency has to be taken more seriously in this case. It's not that we seriously
doubt that TNT could be inconsistent but there is a little doubt, a flicker, a glimmer of a
doubt in our minds, and a proof would help to dispel that doubt.
        But what means of proof would we like to see used? Once again, faced with the
recurrent question of circularity. If we use all the equipment in a proof about our system
as we have inserted into it, what will we have accomplished? If we could manage to
convince ourselves consistency of TNT, but by using a weaker system of reasoning than
we will have beaten the circularity objection! Think of the way a heavy rope is passed
between ships (or so I read when I was a kid): first a light arrow is fired across the gap,
pulling behind it a thin rope. Once a connection has been established between the two
ships this way, then the heavy rope pulled across the gap. If we can use a "light" system
to show that a system is consistent, then we shall have really accomplished something.
        Now on first sight one might think there is a thin rope. Our goal is to prove that
TNT has a certain typographical property (consistency): that no theorems of the form x
and .~x ever occur. This is similar to trying to show that MU is not a theorem of the
MIU-system. Both are statements about typographical properties of symbol-
manipulation systems. The visions of a thin rope are based on the presumption that facts
about number theory won’t be




Typographical Number Theory                                                             237

needed in proving that such a typographical property holds. In other words, if properties
of integers are not used-or if only a few extremely simple ones are used-then we could
achieve the goal of proving TNT consistent, by using means which are weaker than its
own internal modes of reasoning.
         This is the hope which was held by an important school of mathematicians and
logicians in the early part of this century, led by David Hilbert. The goal was to prove the
consistency of formalizations of number theory similar to TNT by employing a very
restricted set of principles of reasoning called "finitistic" methods of reasoning. These
would be the thin rope. Included among finitistic methods are all of propositional
reasoning, as embodied in the Propositional Calculus, and additionally some kinds of
numerical reasoning. But Gödel’s work showed that any effort to pull the heavy rope of
TNT's consistency across the gap by using the thin rope of finitistic methods is doomed
to failure. Gödel showed that in order to pull the heavy rope across the gap, you can't use
a lighter rope; there just isn't a strong enough one. Less metaphorically, we can say: Any
system that is strong enough to prove TNT's consistency is at least as strong as TNT
itself. And so circularity is inevitable.




Typographical Number Theory                                                             238

